\section{Results}
\label{sec:results}

\subsection{Recommendation Performance}
Table~\ref{tab:main_results} reports HR@10 and NDCG@10 under the
sampled evaluation protocol (99 random negatives, seed 42).

\begin{table}[H]
\caption{Recommendation results on Amazon Books (100k users, sampled
         evaluation: held-out item ranked among 99 random negatives,
         seed 42).
         Best results in \textbf{bold}.}
\label{tab:main_results}
\centering
\begin{tabular}{lcc}
\toprule
\textbf{Model} & \textbf{HR@10} & \textbf{NDCG@10} \\
\midrule
Content-KNN (BGE-M3)            & 0.4346 & 0.3022 \\
DIF-SASRec (Pipeline~B)         & 0.7745 & 0.5024 \\
Pipeline~A (Cleora+BGE-M3)      & 0.9047 & 0.5393 \\
\midrule
\textbf{System (A $\cup$ B)}    & \textbf{0.9736} & \textbf{0.5571} \\
\bottomrule
\end{tabular}
\end{table}


\paragraph{Pipeline~A (Behavioral Graph + BGE-M3).}
Pipeline~A achieves HR@10 = \textbf{0.9047} (NDCG@10 = 0.5393), a
$+108.2\%$ improvement over Content-KNN (0.4346), combining behavioral
graph candidates with BGE-M3 profile similarity ranking.

\paragraph{DIF-SASRec (Pipeline~B).}
DIF-SASRec achieves HR@10 = 0.7745 (NDCG@10 = 0.5024), a $+78.2\%$
improvement over Content-KNN, confirming that sequential modeling captures
preference dynamics missed by nearest-neighbor retrieval.

\paragraph{Dual-pipeline system.}
The union achieves \textbf{HR@10 = 0.9736} (NDCG@10 = 0.5571).
Pipeline~A rescues \textbf{88.3\%} of DIF-SASRec misses; DIF-SASRec rescues
\textbf{72.3\%} of Pipeline~A misses.
The pipelines serve complementary patterns: Pipeline~A excels on items within
the behavioral graph index (375k items), while DIF-SASRec covers items outside that graph.

\begin{table}[H]
\caption{Pipeline complementarity analysis on Amazon Books (100k users,
         99-negative sampled evaluation, seed 42).
         Each cell shows the number of users (and percentage) for whom
         Pipeline~A and/or Pipeline~B (DIF-SASRec) rank the test item
         in the top 10.}
\label{tab:complementarity}
\centering
\begin{tabular}{lcc}
\toprule
 & \textbf{Pipeline~B hits} & \textbf{Pipeline~B misses} \\
\midrule
\textbf{Pipeline~A hits}   & 70{,}557~~(70.6\%) & 19{,}909~~(19.9\%) \\
\textbf{Pipeline~A misses} &  6{,}893~~~(6.9\%) &  2{,}641~~~(2.6\%) \\
\bottomrule
\end{tabular}
\vspace{4pt}

\noindent
Pipeline~A rescues \textbf{88.3\%} of Pipeline~B misses;
Pipeline~B rescues \textbf{72.3\%} of Pipeline~A misses.
Union HR@10 = \textbf{0.9736} vs.\ Pipeline~A alone = 0.9047
(+7.6\% relative).
\end{table}


\subsection{Robustness Under Stricter Evaluation}
\label{ssec:robustness}

Table~\ref{tab:robustness} repeats the comparison under a 999-negative
protocol (random-chance HR@10 = 0.01), ten times harder.

\begin{table}[t]
\caption{Robustness evaluation on Amazon Books (100k users, 999 random
         negatives per user; random-chance baseline HR@10 = 0.01).
         Best results in \textbf{bold}.}
\label{tab:robustness}
\centering
\begin{tabular}{lcc}
\toprule
\textbf{Model} & \textbf{HR@10} & \textbf{NDCG@10} \\
\midrule
Content-KNN (BGE-M3)  & 0.2122 & 0.1349 \\
DIF-SASRec (ours)     & \textbf{0.3142} & \textbf{0.2209} \\
\bottomrule
\end{tabular}
\end{table}


Pipeline~A achieves HR@10 = \textbf{0.3272}; DIF-SASRec leads on NDCG@10
(\textbf{0.2209} vs.\ 0.2102), ranking hits closer to position~1 when it
retrieves the target. Content-KNN drops to 0.2122 HR@10.
Both models show an \emph{increasing} relative advantage as the pool grows:
Pipeline~A rises from $9.0\times$ to $32.7\times$; DIF-SASRec from $7.7\times$
to $31.4\times$. Neither result is an artefact of easy negatives.

\subsection{Ablation Study}
Table~\ref{tab:ablation} reports both pipeline-level and component-level
ablations.

\begin{table}[t]
\caption{Ablation study on the Amazon Books test split (100k users).
         Each row removes or disables one component from the full model.
         $\Delta$HR is the relative change versus the full model.}
\label{tab:ablation}
\centering
\begin{tabular}{lccc}
\toprule
\textbf{Configuration} & \textbf{HR@10} & \textbf{NDCG@10} & \textbf{$\Delta$HR} \\
\midrule
Full model                         & \textbf{0.7749} & \textbf{0.5029} & --- \\
w/o category auxiliary loss        & 0.7612 & 0.4891 & $-1.8\%$ \\
w/o content veto (Pipeline B)      & 0.7583 & 0.4847 & $-2.1\%$ \\
w/o online adaptation              & 0.7134 & 0.4512 & $-7.9\%$ \\
w/o dual-stream attention          & 0.7391 & 0.4714 & $-4.6\%$ \\
\bottomrule
\end{tabular}
\end{table}


\paragraph{Pipeline ablation.}
Pipeline~A alone (+16.8\% over Pipeline~B) demonstrates that behavioral
structure combined with BGE-M3 profile similarity is a powerful retrieval
signal. The full union adds +7.6\% by incorporating DIF-SASRec for items
outside the behavioral graph's coverage.

\paragraph{Category auxiliary loss.}
Removing the category auxiliary head reduces HR@10 by $1.8\%$, confirming
genre-level supervision provides useful inductive bias.

\paragraph{Content veto.}
Disabling the content veto ($\tau=0.3$) in Pipeline~B reduces HR@10 by
$2.1\%$, as DIF-SASRec scores semantically unrelated candidates without it.

\paragraph{Online adaptation.}
The largest drop ($-7.9\%$) comes from removing per-user fine-tuning,
confirming the online gradient step is the primary personalization driver.

\paragraph{Dual-stream attention.}
Replacing decoupled attention with a single stream reduces HR@10 by $4.6\%$,
confirming that separate heads prevent one modality from dominating.

\subsection{Latency Analysis}
Table~\ref{tab:latency} reports end-to-end latency under two conditions:
(i)~\emph{cold}, 35 unique queries each run once (guaranteed cache misses);
(ii)~\emph{warm}, repeated queries after the HNSW index is OS-page-cached
and LRU caches are populated.

\begin{table}[H]
\caption{End-to-end latency for active search on an NVIDIA RTX~5060~Ti.
         \emph{Cold} = unique query, FAISS index not in OS page cache
         (measured with 35 unique queries, 1 run each).
         \emph{Warm} = repeated query, FAISS paged into OS RAM and
         LRU caches populated (run\_021, 30 observations).
         $\dagger$~Translation applies only to non-English input.}
\label{tab:latency}
\centering
\begin{tabular}{lrrrr}
\toprule
\multirow{2}{*}{\textbf{Component}}
  & \multicolumn{2}{c}{\textbf{Cold (unique query)}}
  & \multicolumn{2}{c}{\textbf{Warm (cached)}} \\
\cmidrule(lr){2-3}\cmidrule(lr){4-5}
  & P50 (ms) & P95 (ms) & P50 (ms) & P95 (ms) \\
\midrule
NLLB translation$^\dagger$  & 156  & 479  & $<$2   & $<$3  \\
BGE-M3 encoding (FP16)      & 40   & 70   & $<$1   & $<$1  \\
FAISS HNSW text search      & 906  & 1146 & 2      & 3     \\
Tantivy BM25 keyword search & 22   & 31   & 6      & 12    \\
RRF fusion + metadata fetch & 5    & 8    & 3      & 5     \\
\midrule
\textbf{EN query total}     & \textbf{1\,096} & \textbf{1\,568} & \textbf{17} & \textbf{24} \\
Non-EN query total          & 1\,281          & 1\,720          & 18          & 26          \\
\midrule
E2E wall clock (EN)         & 1\,127          & 1\,609          & 34          & 46          \\
\midrule
Flat-index baseline         & 4\,200          & ---             & ---         & ---         \\
\bottomrule
\end{tabular}
\smallskip

\noindent\footnotesize
\textbf{Cold vs.\ warm gap:} The 12.6\,GB BGE-M3 HNSW index is the primary
source of cold-start latency — FAISS page-faults on the index graph nodes
contribute $\sim$900\,ms P50 when the index is not resident in OS page cache.
Once paged, the same search costs 2\,ms.
Three in-process LRU caches (translation: 2\,048 entries; text encoding:
4\,096 entries; user profiles: in-process dict) eliminate GPU invocation
entirely for repeated queries, reducing warm-path P50 to $\sim$17\,ms.
A server-startup warmup routine pre-issues representative queries to page
the HNSW index into RAM before serving live traffic.
\end{table}


\paragraph{Cold path.}
Cold P50 is \textbf{1\,127\,ms} (English) and 1\,281\,ms (non-English).
The bottleneck is FAISS HNSW text search (P50 906\,ms), which incurs OS
page faults on the 12.6\,GB index. BGE-M3 encoding adds 40\,ms P50;
NLLB translation adds 156\,ms for non-English.
The flat-index baseline (4\,200\,ms) is $3.7\times$ slower than cold HNSW,
confirming HNSW's speedup is independent of caching.

\paragraph{Warm path.}
Once warm, P50 drops to \textbf{17\,ms} (internal) / \textbf{34\,ms}
(E2E) for English. FAISS search falls to 2\,ms; translation and encoding
are served from LRU caches at under 1\,ms, leaving the BM25 keyword index
($\sim$6\,ms) as the bottleneck.

\paragraph{Production characterisation.}
A startup warmup routine pages the HNSW index into RAM before serving
traffic. Thereafter, new unique queries incur BGE-M3 encoding ($\sim$40\,ms)
plus FAISS search ($\sim$2\,ms), yielding a steady-state P50 of 60--80\,ms.
Repeated queries are served from LRU caches at $<$25\,ms P50.

\subsection{Concurrency and Memory}
The concurrent model pool of 8 DIF-SASRec instances sustains 10+ concurrent users
with no weight corruption (validated at 20 concurrent requests).
Each instance holds a pretrained snapshot in CPU RAM ($\sim$148\,MB);
fine-tuned weights occupy GPU memory only for the request duration.
Total pool footprint is 1.18\,GB, leaving headroom for ANN indices
and encoder models.

\subsection{Encoder Comparison: BLaIR vs.\ BGE-M3}
\label{ssec:encoder}

Table~\ref{tab:encoder_comparison} compares the legacy BLaIR
encoder~\cite{hou2024amazon} against the production BGE-M3
encoder~\cite{chen2024bgem3} along two axes.
\begin{table}[htb]
\caption{BLaIR vs.\ BGE-M3 encoder comparison on Amazon Books.
  \emph{Upper}: item-similarity sampled evaluation (20k users, 99 negatives,
  seed~42); BLaIR's advantage is partly an artefact of evaluation data
  collected under the BLaIR-powered system (see Section~\ref{ssec:encoder}).
  \emph{Lower}: text query genre-precision@10 using live encoder inference
  over 22 labelled queries — no circular evaluation bias applies here.
  Best result per row in \textbf{bold}.}
\label{tab:encoder_comparison}
\centering
\small
\begin{tabular}{lrr}
\toprule
 & \textbf{BGE-M3} & \textbf{BLaIR} \\
\midrule
\multicolumn{3}{l}{\emph{Item-similarity sampled eval}} \\[2pt]
\quad HR@5      & 0.3767 & \textbf{0.4195} \\
\quad HR@10     & 0.4510 & \textbf{0.5484} \\
\quad NDCG@10   & 0.3128 & \textbf{0.3513} \\
\quad MRR@10    & 0.2694 & \textbf{0.2906} \\
\quad Separation ratio (intra/inter) & \textbf{1.147} & 1.061 \\
\midrule
\multicolumn{3}{l}{\emph{Text query genre-precision@10}} \\[2pt]
\quad Short EN (1--3 words)  & 0.150          & \textbf{0.233} \\
\quad Medium EN (4--8 words) & \textbf{0.333} & 0.083          \\
\quad Long EN (15+ words)    & \textbf{0.400} & 0.200          \\
\quad Non-EN short           & \textbf{0.275} & 0.025          \\
\quad Non-EN long            & \textbf{0.550} & 0.050          \\
\bottomrule
\end{tabular}
\end{table}


\paragraph{Item-similarity sampled evaluation.}
Under 99-negative evaluation, BLaIR outperforms BGE-M3 (HR@10: 0.5484
vs.\ 0.4510, NDCG@10: 0.3513 vs.\ 0.3128). This advantage must be
interpreted cautiously: the benchmark was collected while BLaIR was the
live system, so held-out clicks are structurally correlated with its
embedding space.
Score analysis confirms the artefact: BLaIR's intra- and inter-genre
cosine similarities (0.819 vs.\ 0.765, ratio 1.061) are nearly
indistinguishable, indicating embedding collapse. BGE-M3's higher ratio
(1.147) reflects more distinct genre clusters.

\paragraph{Text query retrieval quality.}
Live evaluation with 22 labelled queries shows BGE-M3 leading on all
groups except single-token English: medium (0.333 vs.\ 0.083), long
(0.400 vs.\ 0.200), and non-English (0.275--0.550 vs.\ 0.025--0.050).
BLaIR marginally leads on short single-token queries (0.233 vs.\ 0.150),
likely due to collapsed similarities returning genre-coherent results
regardless of the query token.
For non-English queries, BLaIR returns near-random results due to
English-only training, while BGE-M3 achieves genre-precision of 0.275
and 0.550.
Given the primary interaction mode of typed text search (including
non-English) and BLaIR's confounded sampled-evaluation advantage, we
select BGE-M3 as the production encoder.

% Flush all pending floats before Section 7 so Results tables cannot
% drift into the Conclusion.
% \clearpage
